\cotefile{15}
\scanpage{68}
$C$ catégorie de Deuring. On a alors la gr-catégorie $\underline{\mathrm{Equ}}(C)$, qui opère sur $C$, donc

\[
\boxed{\mathrm{Brd}(C)=\pi_0(\underline{\mathrm{Equ}})\quad
\deleted{\text{classes d'isom de }C}\ \text{opère sur Cl. d'isom des }C}
\]

Ceci dit, le quotient est l'ens des classes d'isomorphie des $A$-Algèbres $\unsure{o|\sigma}$ apparentées à $C$ (donc, dans le cas des \unsure{$K$|$\mathrm R$} ordres maximaux\ldots, \inserted{le quotient est formé, \unsure{$\exists$}} des classes de conjugaison de tous ordres maximaux, tandis que $\pi_0(\underline{\mathrm{Equ}})$ correspond \deleted{aux} \inserted{est le groupe des} classes d'isomorphie de $o$-bimodules $=$ \unsure{bi-$o$-réseaux} de $o$ mod aux \unsure{princ.|puiss.|provenant} de \deleted{\unsure{$K$|Pic}} $A^*$\unsure{\}|?}, \deleted{\unsure{Noter|Notons}} et $\mathrm{Iso}(C)$ est l'ens des classes d'isomorphie de $o$-modules : \unsure{droites|droits} ``inv.'' \deleted{de} ou encore des $o$-réseaux \deleted{\unsure{x}} à droite]. Dans le cas \inserted{des} \unsure{$A$ centraux|centraux|$o$ centraux} \deleted{\illegible{2}} \inserted{sur $A$} [on dit $C$ centrée sur $A$) on a de plus un foncteur \uline{\unsure{pleint} fidèle}

\[
\begin{aligned}
\underline{\underline{\mathrm{Pic}}}(A)&\longrightarrow\underline{\mathrm{Equ}}(C)\\
L&\longmapsto(X\longmapsto L\otimes_A X)
\end{aligned}
\]

d'où une suite exacte

\[
0\unsure{\mathrel{-}|\longrightarrow}\mathrm{Pic}(A)\longrightarrow
\deleted{\unsure{\mathrm{Br}}}\mathrm{Brdt}(C)\longrightarrow
\frac{\mathrm{Brdt}(C)}{\mathrm{Pic}(A)}\longrightarrow 0
\]

Dans le cas de Brauer (et ordres maximaux, $o$)

\[
\frac{\mathrm{Brd}\,C}{\mathrm{Pic}(A)}\simeq
\frac{\mathrm{Div}(C)}{\mathrm{Div}(A)}=
\underset{\substack{p\in\mathrm{Max}\,A\\p\ \text{ramifié en }C}}{\coprod}
\mathbb{Z}/\underset{\text{ramification en }p.}{e_p}\mathbb Z
\]
% archive sheet number: 67

